\section{Work-conjugate port constraints and Schur extraction}
\label{sec:port-schur}

The electrical port definition must be independent of how an implementation
chooses impressed-current right-hand sides. vNext therefore defines ports
variationally through exact current constraints and their work-conjugate
Lagrange multipliers.

\subsection{Peak-phasor convention}

The theory fixes
\[
  f_{\rm phys}(t)
  =
  \Rea\{\widetilde f\,e^{+j\omega t}\},
\]
where $\widetilde f$ is a peak phasor. For port peak phasors,
time-averaged electrical power entering the electromagnetic subsystem is
\[
  P_{\rm in}
  =
  \frac12\Rea\{v\adj i\}.
  \label{eq:peak-power}
\]
All loss and heat formulas in the theory use this convention unless explicitly
stated otherwise. Under $e^{+j\omega t}$, passive dielectric loss corresponds
to the sign convention used consistently by the constitutive implementation;
material passivity is checked from power, not inferred from an isolated sign of
an imaginary coefficient.

\subsection{Constrained field solve}

Let $x$ contain all electromagnetic internal unknowns after gauge constraints,
and let
\[
  \mathcal I x=i
\]
be the exact vector of oriented terminal-current constraints. For zero external
field, define the constrained system
\[
 \begin{bmatrix}
   A & -\mathcal I\adj\\
   \mathcal I & 0
 \end{bmatrix}
 \begin{bmatrix}
   x\\v
 \end{bmatrix}
 =
 \begin{bmatrix}
   0\\i
 \end{bmatrix}.
 \label{eq:port-kkt}
\]
The multiplier $v$ is, by definition, the work-conjugate port-voltage vector.
The minus sign in Equation~\eqref{eq:port-kkt} fixes the orientation convention.
If an implementation uses the opposite terminal-current orientation, both
$\mathcal I$ and the corresponding voltage orientation transform together.

Eliminating $x$ gives
\[
  x=A^{-1}\mathcal I\adj v,
\]
\[
  i
  =
  \mathcal I A^{-1}\mathcal I\adj v.
\]
Hence the port admittance and impedance are
\[
  Y
  =
  \mathcal I A^{-1}\mathcal I\adj,
  \qquad
  Z=Y^{-1},
  \label{eq:port-schur}
\]
on the independent port subspace.

\subsection{Current-driven implementation}

For a prescribed current vector, an implementation may solve the KKT system
directly. It may alternatively construct a compatible lifting $x_i$ satisfying
\[
  \mathcal I x_i=i
\]
and solve only for a zero-port-current correction. These formulations are
equivalent only if the lifting preserves the same charge-continuity and gauge
constraints.

\subsection{Voltage-driven implementation}

For prescribed $v$, solve
\[
  Ax=\mathcal I\adj v
\]
and recover
\[
  i=\mathcal I x.
\]
This form is especially convenient for extracting $Y$. Conversion to $Z$
requires a nonsingular independent-port admittance. Floating/common-mode
electrical degrees of freedom are removed according to the circuit/port
topology rather than by arbitrary matrix regularization.

\subsection{Power identity}

Multiplying the field equation by $x\adj$ gives
\[
  x\adj A x
  =
  x\adj\mathcal I\adj v
  =
  i\adj v.
\]
Therefore
\[
  P_{\rm in}
  =
  \frac12\Rea\{x\adj A x\}
  =
  \frac12\Rea\{v\adj i\}.
\]
The real part of the physical operator decomposes into conductor, dielectric,
magnetic, and outward channels. This is the primary port-to-field power closure
identity used by REFERENCE and CERTIFIED modes.

\subsection{Reciprocity}

If the reciprocal physical operator and port functionals satisfy the
corresponding complex-symmetric reciprocity relation, then
\[
  Y^\T=Y,\qquad Z^\T=Z.
\]
Reciprocity is therefore a property of the operator and work-conjugate port
definition, not a matrix symmetrization step applied after solving.

\subsection{Multi-terminal networks}

For $N_t$ physical terminals, let $C_p$ map terminal currents to $P$ independent
port currents while enforcing Kirchhoff conservation:
\[
  i_t=C_p i,\qquad \mathbf 1^\T i_t=0
\]
within each isolated electrical network. The port current functional is
\[
  \mathcal I=C_p^\T\mathcal I_t
\]
under the chosen orientation. Port voltages use the dual terminal-potential
combination. This permits differential, multi-winding, and future branched
conductor networks without changing the field operator.

\subsection{Contact impedance}

A physical terminal/contact impedance is represented explicitly by
\[
  Z_{\rm port}
  =
  Z_{\rm field}+Z_{\rm contact},
\]
where
\[
 \Herm Z_{\rm contact}\PSD
\]
for a passive contact. Contact terms are never used as numerical regularizers
for an ill-conditioned field solve.

\subsection{Teacher outputs}

For each scene/frequency the conductor teacher stores:
\[
  \{Y,Z,X_{\rm port},D_{\rm channels},
    \eta_{\rm solve},\eta_{\rm power}\},
\]
plus modal field coefficients needed for continuous loss queries. The KKT
multiplier or equivalent Schur extraction is the canonical definition against
which alternative source formulations are regression-tested.
